\documentclass[11.9pt,letterpaper]{article}
\usepackage[utf8]{inputenc}
\usepackage[spanish]{babel}
\usepackage{amsmath}
\usepackage{amsfonts}
\usepackage{diagbox}
\usepackage{amssymb}
\usepackage{makeidx}
\usepackage{array}
\usepackage{multirow}
\usepackage{multicol}
\usepackage{booktabs}
\usepackage{colortbl}
\usepackage{setspace}
\usepackage{graphicx}
\usepackage{float}
\usepackage[xcdraw]{xcolor}
\usepackage{colortbl}
\usepackage[margin=1.5cm]{geometry}
\definecolor{TAREA}{HTML}{FF65D4}
\definecolor{EXAMEN}{HTML}{99CE3E}
\definecolor{VIRTUAL}{HTML}{5aa3ec}
\title{
\vspace*{-1.5cm}
\Large \textbf{Universidad Nacional Autónoma de Honduras}\\[0.5cm]
\includegraphics[scale=0.4]{UNAH_Gris}\\[0.1cm]
\Large Escuela de Matemática y Ciencias de la Computación\\
\Large Departamento de Matemática Aplicada\\
\large Análisis Numérico, MMM412\\
\large Tercer periodo académico 2024.
}
\begin{document}
\maketitle
\begin{multicols}{2}
\renewcommand{\arraystretch}{1.2}
\begin{center}
\begin{tabular}{c}
\cellcolor{gray!40}{GENERALIDADES DEL CURSO}
\end{tabular}
\end{center}
Esta es la información general del curso.
\begin{itemize}
\item \textbf{Asignatura:} Análisis Numérico.
\item \textbf{Créditos:} 3.
\item \textbf{Docente:} Myrian Gonzalez.
\item \textbf{Email:} myrian.gonzalez@unah.edu.hn
\end{itemize}
\begin{center}
\begin{tabular}{c}
\cellcolor{gray!40}{METODOLOGÍA DE ENSEÑANZA}
\end{tabular}
\end{center}
\hspace*{0.5cm}El proceso de enseñanza-aprendizaje de este espacio de aprendizaje se desarrollará mediante la modalidad presencial. Las actividades se han desarrollado para que el estudiante por medio de diversas prácticas alcance los conocimientos del curso.
\begin{center}
\begin{tabular}{c}
\cellcolor{gray!40}{MATERIALES Y RECURSOS DIDÁCTICOS}
\end{tabular}
\end{center}
\hspace*{0.5cm} Para el proceso de enseñanza se utilizarán los siguientes recursos. 
\begin{itemize}
\item Proyector: Se usará un conjunto de presentaciones como apoyo al estudiante para el desarrollo y comprensión de la clase.
\item Campus Virtual. En este espacio se asignarán diferentes actividades relacionadas con la clase. 
\item Software: Se utilizará un lenguaje de programación de álgebra computacional para el desarrollo y demostración de los temas de la clase.
\end{itemize}
\begin{center}
\begin{tabular}{c}
\cellcolor{gray!40}{CONTENIDOS}
\end{tabular}
\end{center}
\hspace*{0.5cm} Este curso esta diseñado para que el estudiante controle diferentes aspectos técnicos y prácticos referentes al análisis numérico:
\begin{itemize}
\item Preliminares del cálculo.
\item Métodos para la busqueda de raíces.
\item Métodos para la aproximación de funciones.
\item Integración Numérica.
\item Diferenciación Numérica.
\item Métodos para la resolución de problemas de valor inicial. 
\item Problemas con valores en la frontera para ecuaciones diferenciales ordinarias.
\end{itemize} 
%==========================================================================
\begin{center}
\begin{tabular}{c}
\cellcolor{gray!40}{EVALUACIÓN}
\end{tabular}
\end{center}
\vspace*{0.5cm} Para la evalución de los contenidos aprendidos en la clase se considerarán las siguientes actividades:
\begin{table}[H]
\centering
\begin{tabular}{|l|l|c|}
\hline
Actividad  & Horario & Porcentaje\\\hline\hline
Prueba I&11:00 $-$12:00 &9\%\\
Prueba II&11:00 $-$12:00 &14\%\\
Prueba III&11:00 $-$12:00 &14\%\\
Prueba IV&11:00 $-$12:00 &14\%\\
Prueba V&11:00 $-$12:00 &14\%\\
Prueba VI&11:00 $-$12:00 &14\%\\
Tarea I&11:00 &4\%\\
Tarea II&11:00 &4\%\\
Tarea III&11:00 &4\%\\
Prueba Virtual I &2 días &3\%\\
Prueba Virtual II &2 días &3\%\\
Prueba Virtual III &2 días&3\%\\\hline
\end{tabular}
\end{table}
%==========================================================================
\newpage
\hspace*{0.5cm} 
\begin{center}
\begin{tabular}{c}
\cellcolor{gray!40}{PROGRAMACIÓN}
\end{tabular}
\end{center}
\begin{center}
\begin{tabular}{ccccc|ccccc}
\multicolumn{5}{c|}{\cellcolor{red!40}{SEPTIEMBRE}}&\multicolumn{5}{c}{\cellcolor{blue!40}{OCTUBRE}}\\
S.&L&M&M&J&S.&L&M&M&J\\
I &\cellcolor{gray!40}{9}&\cellcolor{gray!40}{10}&11&12&IV&7&8&9&\cellcolor{white}{10}\\
II&\cellcolor{gray!40}{16}&\cellcolor{gray!40}{17}&18&\cellcolor{white}{19}&\cellcolor{VIRTUAL}{V}&14&\cellcolor{gray!40}{15}&\cellcolor{EXAMEN}{$16\textcolor{TAREA}{\blacksquare}$}&\cellcolor{white}{17}\\
III&\cellcolor{white}{23}&24&\cellcolor{EXAMEN}{25}&\cellcolor{white}{26}&VI&\cellcolor{white}{21}&\cellcolor{white}{22}&\cellcolor{white}{23}&\cellcolor{white}{24}\\
&&&&&VII&\cellcolor{white}{28}&29&\cellcolor{EXAMEN}{30}&31\\
\multicolumn{5}{c|}{\cellcolor{cyan!40}{NOVIEMBRE}}&\multicolumn{5}{c}{\cellcolor{green!40}{DICIEMBRE}}\\
S.&L&M&M&J&S.&L&M&M&J\\
VIII&\cellcolor{white}{4}&5&6&7&\cellcolor{white}{XII}&\cellcolor{white}{2}&3&4&5\\
\cellcolor{VIRTUAL}{IX}&\cellcolor{white}{11}&12&\cellcolor{EXAMEN}{$13\textcolor{TAREA}{\blacksquare}$}&\cellcolor{white}{14}&\cellcolor{VIRTUAL}{XIII}&\cellcolor{white}{9}&10&\cellcolor{EXAMEN}{$11\textcolor{TAREA}{\blacksquare}$}&12\\
X&18&19&20&\cellcolor{white}{21}&XIV&\cellcolor{gray!40}{16}&\cellcolor{gray!40}{17}&\cellcolor{gray!40}{18}&\\
XI&\cellcolor{white}{25}&26&\cellcolor{EXAMEN}{27}&28&&&&&\\
\end{tabular} 
\end{center}
\begin{center}
\begin{tabular}{lp{0.41\textwidth}}
\cellcolor{TAREA}{\ }&\textbf{Tarea}: Este color representa las fechas en que se entregarán las tareas.\\
\cellcolor{VIRTUAL}{\ }&\textbf{Prueba virtual}: Este color indica la semana en que se realizará una prueba en el campus virtual. Dichas pruebas estarán abiertas durante todo el fin de semana; por ejemplo en la semana V habrá una prueba en las fechas 19 y 20 de octubre. \\
\cellcolor{EXAMEN}{\ }&\textbf{Prueba presencial}: Este color representa las fechas en que se realizará una prueba en el aula de clases.\\
\cellcolor{gray!40}{\ }&Estas fechas corresponden a días feriados o ingreso de calificaciones finales del periodo.
\end{tabular}
\end{center}
\begin{center}
\begin{tabular}{c}
\cellcolor{gray!40}{BIBLIOGRAFÍAS}
\end{tabular}
\end{center}
\begin{itemize}
\item Burden, R. L. (2017). Análisis numérico.
\item Sauer, T., \& Murrieta, J. E. M. (2013). Análisis numérico. México: Pearson Educación.
\end{itemize}
\end{multicols}
\begin{center}
\end{center}
\end{document}
